\documentclass[12pt,a4paper]{article}

\usepackage[utf8]{inputenc}
\usepackage[T1]{fontenc}
\usepackage[margin=2.5cm]{geometry}
\usepackage{hyperref}
\usepackage{microtype}
\usepackage{parskip}
\usepackage{lmodern}
\usepackage{amsmath}

\hypersetup{colorlinks=true, urlcolor=blue!60!black}

\pagestyle{empty}

\begin{document}

\begin{flushright}
Kundai Farai Sachikonye \\
Auf der Lichtung 19, 82178 Puchheim, Germany \\
\href{https://github.com/fullscreen-triangle}{github.com/fullscreen-triangle} \\
\texttt{kundai.sachikonye@bitspark.com} \\
(+49) 1626386344 \\[6pt]
1 June 2026
\end{flushright}

\vspace{1em}

Dr.\ Chris Carnie \\
Klinik Kinderheilkunde III (Onkologie, H\"amatologie und Immunologie, Pneumologie) \\
Hopp Children's Cancer Center Heidelberg (KiTZ) / Universit\"atsklinikum Heidelberg \\
Im Neuenheimer Feld 350, 69120 Heidelberg \\
Job-ID: V000015650

\vspace{1.5em}

\textbf{Re: PhD position --- Genome Stability and Nucleotide Metabolism in Cancer
(V000015650)}

\vspace{1em}

Dear Dr.\ Carnie,

I am applying for the doctoral position in your group on the Emmy
Noether-funded project ``Nucleotide Metabolism and Genome Stability in
Cancer,'' and I am grateful to you for drawing my attention to it. The project
sits very close to where I trained: my research began on chemically-induced DNA
damage, and my work since has been computational. Because the position spans
wet lab, dry lab, and bioinformatic analysis of the ’omics data the projects
generate, I am applying as someone who can contribute across that whole span
rather than at one end of it.

\textbf{The biology is where I began.}
My bachelor's thesis was an experimental study of DNA interstrand crosslinking
--- the formation of crosslinks and the cellular response to crosslinking
agents. That is directly adjacent to the antimetabolite biology at the centre of
your project: both crosslinkers and antimetabolite drugs are agents that perturb
the cell and impose potentially cytotoxic DNA damage, studied through the same
repertoire of cell-survival, repair, and damage-response assays. I have a
molecular and cell-biology grounding (BSc in Biochemistry and Cell Biology, MSc
in Pharmaceutical Biotechnology), bench training including mammalian cell
culture and microscopy, and a genuine wish to return to laboratory work in
exactly this area --- the interplay between a metabolic perturbation, the DNA
damage it causes, and the pathways that determine whether a cancer cell
survives it.

\textbf{The analysis is my strongest contribution.}
The position lists bioinformatic aptitude with R or Python as desirable, and the
tasks include dry-lab investigation and analysis of the ’omics data you
generate. This is where I am strongest. My doctoral work was at TUM's Chair of
Proteomics and Bioanalytics, so mass-spectrometry and proteomic / metabolomic
data analysis is home ground for me --- directly relevant to the metabolomics in
your toolkit. I have built a multi-omics inference pipeline integrating
genotype, transcriptomics, and metagenomics, validated against 1000~Genomes,
gnomAD, and UK~Biobank. Genome-wide CRISPR screens, at the analysis layer, are a
genetic-interaction problem --- identifying the dependencies that emerge only
against a metabolic or repair-deficient background --- and I would be glad to
own that analysis for the group while contributing at the bench. A candidate who
can both generate the data and analyse it is, I think, what makes the dual
wet/dry structure of these projects work.

\textbf{Relevant computational work.}
My independent research has been in computational oncology, drug metabolism, and
pharmacology: tumour characterisation, neoantigen and MHC-binding prediction
(AUC~0.81 on 2{,}847 IEDB peptides), and pharmacokinetic / pharmacodynamic and
drug-metabolism modelling validated against NIST reference data --- the last of
these directly concerned with how antimetabolite-class compounds are processed
and act. These tools are openly available (linked above) should they be of
interest.

\textbf{Background.}
I hold an MSc in Computational Biology (Universit\"at T\"ubingen), an MSc in
Pharmaceutical Biotechnology (HAW Hamburg), and a BSc in Biochemistry and Cell
Biology (Jacobs University Bremen). I conducted doctoral research in
Computational Lipidomics at TUM (2019--2021) before transitioning to independent
research. I have lived and worked in Germany continuously since 2011; English is
native and my German is conversational, which suits an English-language working
environment well.

I am available at the earliest possible date and would welcome the opportunity
to discuss how this combination of experimental DNA-damage background and
computational analysis could contribute to your team.

\vspace{1.5em}
Yours sincerely,

\vspace{2.5em}
\textbf{Kundai Farai Sachikonye}

\end{document}
